\begin{name}
	{\tenchude}
	{TOÁN 10}
	{LỚP TOÁN THẦY PHÁT}
	{Life isn’t about finding yourself. Life is about creating yourself}
\end{name}
\Opensolutionfile{ansbook}[ans/ansbookDe4]
\TN
\Opensolutionfile{ans}[ans/ansDe4-TN1]
\begin{ex} %[0D8N1-2]
Nam muốn qua nhà Lan để cùng Lan tới trường. Từ nhà Nam tới nhà Lan có $3$ con đường, từ nhà Lan đến trường có $5$ con đường. Hỏi Nam có bao nhiêu cách chọn đường đi từ nhà đến trường?
\choice
{$8$}
{$243$}%[0D8H1-2]
{$10$}
{\True $15$}
\loigiai{
Có $3$ cách đi từ nhà Nam tới nhà Lan.\\
Sau đó đi từ nhà Lan tới trường có $5$ cách.\\
Theo quy tắc nhân ta có $3\cdot 5=15$ cách.
}
\end{ex}

\begin{ex}%[0D8N2-1]%[Dự án đề kiểm tra Toán 10 HK2 NH23-24- Nguyễn Cường]%[Sở GD-ĐT Bắc Ninh]
Khẳng định nào sau đây là đúng?
\choice
{$\mathrm{A}_8^3=24$}
{$\mathrm{A}_8^3=512$}
{\True $\mathrm{A}_8^3=336$}
{$\mathrm{A}_8^3=56$}
\loigiai{
Ta có $\mathrm{A}_8^3=336$.
}
\end{ex}

\begin{ex}%[0D8N3-3]
Nếu bốn số hạng đầu của một hàng trong tam giác Pascal được ghi lại là  $1; 16; 120; 560$. Khi đó bốn số hạng đầu của hàng kế tiếp là
\choice
{$1; 16; 2312; 67200$}
{$1; 17; 2312; 67200$}
{$1; 17; 126;  680$}
{\True $1; 17; 136; 680$}
\loigiai{
Theo quy luật tam giác Pascal ta có
\begin{align*}
&1+16=17 \\
&16+120=136 \\
& 120+560=680.
\end{align*}
Vậy bốn số hạng đầu của hàng kế tiếp là $1; 17; 136; 680$.}
\end{ex}

\begin{ex}%[10-11-GK2-2425]%[Đỗ Minh Phúc]%[0D6H1-1]
	Nếu lấy $3{,}1$ làm giá trị gần đúng của $3{,}14$ thì sai số tuyệt đối là
	\choice
	{$-6{,}24$}
	{$-0{,}04$}
	{$6{,}24$}
	{\True $0{,}04$}
	\loigiai{
	Sai số tuyệt đối là $\Delta_a=\left|3{,}14-3{,}1\right|=0{,}04$.
	}
\end{ex}

\begin{ex}%[HK1, THPT Chuyên Hùng Vương - Phú Thọ, 2023-2024]%[Lê Quốc Hiệp, 10-11EX-HK1-2324]%[0D6H4-2]
	Mẫu số liệu sau cho biết chiều cao (đơn vị cm) của các bạn trong tổ
	\[\begin{array}{llllllll}163 & 159 & 172 & 167 & 165 & 168 & 170 & 161\end{array}\]
	Khoảng biến thiên của mẫu số liệu này bằng
	\choice
	{$11$}
	{$10$}
	{\True $13$}
	{$12$}
	\loigiai
	{
		Sắp xếp mẫu số liệu theo thứ tự không giảm
		\[
			\begin{array}{llllllll}
				159 & 161 & 163 & 165 & 167 & 168 & 170 & 172
			\end{array}
		\]
		Khoảng biến thiên $R=x_8-x_1=172-159=13$.
	}
\end{ex}

\begin{ex}%[soạn để GHK, HK]%[Huỳnh Xuân Tín]%[0D0N1-2]
	Gieo hai đồng tiền một lần. Kí hiệu $S$, $N$ lần lượt để chỉ đồng tiền xuất hiện mặt sấp, đồng tiền xuất hiện mặt ngửa. Mô tả không gian mẫu nào dưới đây là đúng?
	\choice
	{$\Omega=\{S; N\}$}
	{$\Omega=\{NN; SS\}$}
	{$\Omega=\{SN; NS\}$}
	{\True $\Omega=\{SN; NS; SS; NN\}$}
	\loigiai{
		Không gian mẫu $\Omega=\{SN; NS; SS; NN\}$.
	}
\end{ex}

\begin{ex}%[0D0N1-3]%[tex hóa đề CK2 - 2025 - Huyên Nguyễn]
	Tâm đi từ nhà mình đến nhà Huyền, cùng Huyền đi đến nhà Linh chơi. Biết từ nhà Tâm đến nhà Huyền có $5$ con đường đi. Từ nhà Huyền đến nhà Linh có $7$ con đường đi. Hỏi có bao nhiêu cách để Tâm đi đến nhà Linh mà phải đi qua nhà Huyền?
	\choice{$12$}{\True$35$}{$20$}{$25$}
	\loigiai{
		\begin{itemize}
			\item Từ nhà Tâm đến nhà Huyền có $5$ cách.
			\item Từ nhà Huyền đến nhà Linh có $7$ cách.
		\end{itemize}
		Vậy theo quy tắc nhân, Tâm đi đến nhà Linh mà phải đi qua nhà Huyền sẽ có $5\cdot 7=35 $ cách.
	}
\end{ex}

\begin{ex}%[0H9N3-1]
	Vectơ nào dưới đây là một vectơ chỉ phương của đường thẳng song song với trục $Ox$?
	\choice
	{\True $\vec u=\left( 1;0 \right)$}
	{$\vec u=\left((1;-1\right)$}
	{$\vec u=\left(1;1\right)$}
	{$\vec u=\left(0;1\right)$}
	\loigiai{
		Vectơ $\vec i=\left(1;0\right)$ là một vectơ chỉ phương của trục $Ox$.\\
		Các đường thẳng song song với trục $Ox$ có một vectơ chỉ phương là $\vec u=\vec i=\left(1;0\right)$.
	}
\end{ex}

\begin{ex}%[10-11-12EX-GHK2-2425]%[Dương Quang]%[0H9N3-4]
	Trong mặt phẳng $Oxy$, cho hai đường thẳng $\Delta_1\colon a_1x+b_1y+c_1=0$ ($a_1^2+b_1^2>0$) và $\Delta_2\colon a_2x+b_2y+c_2=0$ ($a_2^2+b_2^2>0$). Trong các công thức sau đây, công thức nào đúng?
	\choice
	{$\cos\left(\Delta_1,\Delta_2\right)=\dfrac{\left|a_1b_1+a_2b_2\right|}{\sqrt{a_1^2+a_2^2}\cdot\sqrt{b_1^2+b_2^2}}$}
	{$\cos\left(\Delta_1,\Delta_2\right)=\dfrac{\left|a_1a_2+b_1b_2\right|}{\sqrt{a_1^2+a_2^2}\cdot\sqrt{b_1^2+b_2^2}}$}
	{\True $\cos\left(\Delta_1,\Delta_2\right)=\dfrac{\left|a_1a_2+b_1b_2\right|}{\sqrt{a_1^2+b_1^2}\cdot\sqrt{a_2^2+b_2^2}}$}
	{$\cos\left(\Delta_1,\Delta_2\right)=\dfrac{\left|a_1a_2+b_1b_2\right|}{\sqrt{a_1^2+b_1^2}\cdot\sqrt{a_2^2+b_2^2}}$}
	\loigiai{
		Ta có $\cos\left(\Delta_1,\Delta_2\right)=\dfrac{\left|a_1a_2+b_1b_2\right|}{\sqrt{a_1^2+b_1^2}\cdot\sqrt{a_2^2+b_2^2}}$.
	}
\end{ex}

\begin{ex}%[0H9N3-3]
	Trong mặt phẳng tọa độ $Oxy$, cho hai đường thẳng $d_1\colon 2x+3y+4=0$ và $d_2\colon 4x+6y-1=0$. Hãy chọn khẳng định đúng.
	\choice
	{$d_1$ và $d_2$ cắt nhau}
	{\True $d_1$ và $d_2$ song song với nhau}
	{$d_1$ và $d_2$ vuông góc nhau}
	{$d_1$ và $d_2$ trùng nhau}
	\loigiai{
		Đường thẳng $d_1\colon 2x+3y+4=0$ có vectơ pháp tuyến $\overrightarrow{n}_1=(2;3)$.\\
		Đường thẳng $d_2\colon 4x+6y-1=0$ có vectơ pháp tuyến $\overrightarrow{n}_2=(4;6)$.\\
		Ta thấy $\overrightarrow{n}_2=2\overrightarrow{n}_2$ nên hai đường thẳng song song hoặc trùng nhau.\\
		Mặt khác, ta lại có điểm $M(-2;0)$ thuộc đường thẳng $d_1$ nhưng không thuộc đường thẳng $d_2$ nên hai đường thẳng này không trùng nhau.\\
		Vậy hai đường thẳng $d_1$ và $d_2$ song song với nhau.
	}
\end{ex}

\begin{ex}%[0H9N5-1]%[Dự án đề kiểm tra Toán 10, 11 HKII NH23-24-BCTuan]%[THPT LTV-Hà Nội]
	Trong mặt phẳng tọa độ $Oxy$, cho elip có phương trình  $\dfrac{x^2}{4}+\dfrac{y^2}{1}=1$. Một tiêu điểm của elip có tọa độ là
	\choice
	{$D\left(0;\sqrt{5}\right)$}
	{$B\left(0;\sqrt{3} \right)$}
	{$C\left(\sqrt{5};0\right)$}
	{\True $A\left(\sqrt{3};0\right)$}
	\loigiai{
		Các tiêu điểm của elip có tọa độ là $(c;0)$ và $(-c;0)$.\\
		Ta có $a^2=4$, $b^2=1$ nên $c^2=a^2-b^2=3\Leftrightarrow c=\sqrt{3}$.\\
		Vậy một tiêu điểm của elip có tọa độ là $A\left(\sqrt{3};0\right)$.
	}
\end{ex}

\begin{ex}%[0H9N5-5]%[Dự án đề kiểm tra Toán 10 HKII NH23-24- Đoàn Minh Tâm]%[THPT Mạc Đĩnh Chi]
	Trong mặt phẳng tọa độ $Oxy$, hypebol $(H)$ có độ dài trục thực bằng $ 6$ và tiêu cự bằng $ 10$ có phương trình chính tắc là
	\choice
	{\True $(H)\colon\dfrac{x^2}{9}-\dfrac{y^2}{16}=1$}
	{$(H)\colon\dfrac{x^2}{16}-\dfrac{y^2}{9}=1$}
	{$(H)\colon\dfrac{x^2}{25}-\dfrac{y^2}{9}=1$}
	{$(H)\colon\dfrac{x^2}{36}-\dfrac{y^2}{64}=1$}
	\loigiai{
		Độ dài trục thực bằng $6$ nên $2a=6 \Rightarrow a=3\Rightarrow a^2=9$.\\
		Tiêu cự bằng $10$ nên $2c=10 \Rightarrow c=5$.\\
		Ta có $b^2=c^2-a^2=25-9=16$.\\
		Vậy phương trình chính tắc của hypebol $(H)$ là $\dfrac{x^2}{9}-\dfrac{y^2}{16}=1$.
	}
\end{ex}
\Closesolutionfile{ans}

\TNTF
\Opensolutionfile{ans}[ans/ansDe4-TN2]
\begin{ex}%[0D6H4-2]
	Điểm trung bình các môn trong kỳ thi tốt nghiệp trung học phổ thông năm 2024 được thống kê trong bảng sau:
	\begin{center}
		\begin{tabular}{|c|c|c|c|c|c|c|c|c|c|}
			\hline
			\textbf{Môn học} & Toán & Văn  & Vật lý & Hóa học & Sinh học & Lịch sừ & Địa lý & GDCD & Ngoại ngữ \\
			\hline
			\textbf{Điểm}    & 6,45 & 7,23 & 6,67   & 6,68    & 6,28     & 6,57    & 7,19   & 8,16 & 5,51      \\
			\hline
		\end{tabular}
	\end{center}
	Các khẳng định sau đúng hay sai?
	\choiceTF[t]
	{\True Điểm trung bình của $9$ môn thi tốt nghiệp năm $2024$ (làm tròn đến hàng phần trăm) là $6{,}75$}
	{\True Khoảng biến thiên của mẫu số liệu trên là $2{,}65$}
	{\True Quy tròn số trung bình với độ chính xác $d=0{,}01$ bằng $a=6{,}8$}
	{Sai số tương đối của số gần đúng trên xấp xỉ bằng $1,48\%$}
	\loigiai{
		\begin{itemchoice}
			\itemch Vì $\bar{x}=\dfrac{6,45+7,23+6,67+6,68+6,28+6,57+7,19+8,16+5,51}{9}=\dfrac{3037}{450} \approx 6,75$
			\itemch Giá trị bé nhất của mẫu số liệu trên là: $5{,}51$. \\
			Giá trị lớn nhất của mẫu số liệu trên là: $8{,}16$. \\
			Do đó: khoảng biến thiên của mẫu số liệu trên là $8{,}16-5{,}51=2{,}65$.
			\itemch Quy tròn số trung bình với độ chính xác $d=0{,}01$ bằng $a=6{,}8$.
			\itemch Sai số tương đối của số gần đúng
			$$\delta=\dfrac{\left|6{,}8-6{,}75\right|}{6{,}8}\approx 0{,}74\%$$
		\end{itemchoice}
	}
\end{ex}

\begin{ex}%[0H9H3-2]
	Cho ba điểm $M(1;-2)$, $N(4;-1)$ và $P(1;2)$. 
	\choiceTF
	{\True Đường thẳng $\Delta$ đi qua $M$ và $N$ có vectơ pháp tuyến là $\vec{n}=(1;-3)$}
	{\True Phương trình tổng quát của đường thẳng $\Delta \colon x-3y-7=0$}
	{Đường tròn ngoại tiếp tam giác $MNP$ có tâm là $I(1;0)$}
	{\True Phương trình đường tròn ngoại tiếp tam giác $MNP$ là $(x-2)^2+y^2=5$}
	\loigiai{
		\begin{itemchoice}
			\itemch Đường thẳng $\Delta \colon \heva{&x=1+3t\\&y=-2+t}$ có vectơ chỉ phương là $\vec{u}=\left(3;1\right)$ nên nhận $\vec{n}=\left(1;-3\right)$ là một vectơ pháp tuyến.
			\itemch Đường thẳng $\Delta$ đi qua điểm $M\left(1;-2\right)$ nên phương trình tổng quát của $\Delta$ là \[\left(x-1\right)-3\left(y+2\right)=0\Leftrightarrow x-3y-7=0.\]
			\itemch Đường tròn ngoại tiếp tam giác $MNP$ có tâm $I(x;y)$ nên 
			\[\heva{&IM=IN \\ &IM=IP}\Leftrightarrow \heva{&(x-1)^2+(y+2)^2=(x-4)^2+(y+1)^2 \\ &(x-1)^2+(y+2)^2=(x-1)^2+(y-2)^2}\Leftrightarrow \heva{&x^2+y^2-2x+4y+5=x^2+y^2-8x+2y+17 \\ &x^2+y^2-2x+4y+5=x^2+y^2-2x-4y+5}\Leftrightarrow \heva{&6x+2y=12 \\ &8y=0}\Leftrightarrow \heva{&x=2 \\ &y=0}.\]
			\itemch Đường tròn ngoại tiếp tam giác $MNP$ có tâm $I(2;0)$ và bán kính $R=IM=\sqrt{(2-1)^2+(0+2)^2}=\sqrt{5}$ nên phương trình đường tròn ngoại tiếp tam giác $MNP$ là $(x-2)^2+y^2=5$.
		\end{itemchoice}
	}
\end{ex}
\Closesolutionfile{ans}

\TNSA
\Opensolutionfile{ans}[ans/ansDe4-TN3]
\begin{ex}%[0D8H1-3]
Một cửa hàng đồ chơi có $3$ loại đồ chơi gồm: $8$ kiểu ô tô khác nhau, $7$ kiểu máy bay khác nhau và $10$ kiểu đồ chơi xếp hình khác nhau. Bạn Minh muốn mua hai đồ chơi khác loại. Hỏi bạn Minh có bao nhiêu cách chọn mua?
\shortans{206}
\loigiai{
Ta có ba trường hợp sau
\begin{itemize}
\item Chọn mua $1$ ô tô và $1$ máy bay có $8\cdot 7=56$ cách.
\item Chọn mua $1$ ô tô và $1$ đồ chơi xếp hình có $8\cdot 10=80$ cách.
\item Chọn mua $1$ máy bay và $1$ đồ chơi xếp hình có $7\cdot10=70$ cách.
\end{itemize}
Theo quy tắc cộng có tổng cộng là $56+80+70=206$ cách chọn mua hai đồ chơi khác loại.
}
\end{ex}

\begin{ex}%[0D8V2-4]
Một nhóm $9$ người gồm ba người đàn ông, bốn phụ nữ và hai đứa trẻ đi xem phim. Hỏi có bao nhiêu cách xếp họ ngồi trên một hàng ghế sao cho mỗi đứa trẻ ngồi giữa hai người phụ nữ và không có hai người đàn ông nào cạnh nhau?
% \shortans[oly]{864}
\loigiai{
Số cách xếp $4$ người phụ nữ là $4! =24$.\\
Có $3$ vị trí giữa $4$ người phụ nữ để xếp $2$ đứa trẻ. Do đó có $\mathrm{A}_3^2=6$ cách xếp $2$ đứa trẻ.\\
$3$ vị trí còn lại cho $3$ người đàn ông là đầu hàng ghế, cuối  hàng ghế và vị trí còn trống giữa $3$ người phụ nữ. Do đó có $3!=6$ cách xếp $3$ người đàn ông.\\
Vậy có $24\cdot 6\cdot 6=864$ cách sắp xếp $9$ người thỏa mãn yêu cầu bài toán.
}
\end{ex}

\begin{ex}%[0H9V5-0]%[Dự án đề kiểm tra Toán 11 GHKII, HKII NH23-24- Huỳnh Xuân Tín]%[Sở GH-ĐT Bắc Giang]
\immini{Hình bên là một tấm giấy hình chữ nhật có chiều dài bằng $24 \mathrm{~cm}$ và chiều rộng bằng $16 \mathrm{~cm}$, trên đó có một đường tròn và hai nhánh của một hypebol. Tìm tiêu cự của hypebol (Đơn vị: cm, kết quả làm tròn đến chữ số thập phân thứ nhất).}{\begin{tikzpicture}[scale=1, font=\footnotesize, line join=round, line cap=round, >=stealth]
\def\a{2};
\def\b{1.788854382};
\pgfmathsetmacro\c{sqrt((\a)^2+(\b)^2)};
\draw[name path=h1,smooth, samples=200, domain=\a:3] plot (\x,{((\b)/(\a))*sqrt((\x)^2-(\a)^2)});
\draw[yscale=-1,name path=h2,smooth, samples=200, domain=\a:3] plot (\x,{((\b)/(\a))*sqrt((\x)^2-(\a)^2)});
\draw[xscale=-1,name path=h2,smooth, samples=200, domain=\a:3] plot (\x,{((\b)/(\a))*sqrt((\x)^2-(\a)^2)});
\draw[xscale=-1,yscale=-1,name path=h2,smooth, samples=200, domain=\a:3] plot (\x,{((\b)/(\a))*sqrt((\x)^2-(\a)^2)});
\def\xm{2};
\pgfmathsetmacro\ym{((\b)/(\a))*sqrt((\xm)^2-(\a)^2)};
\path
(0,0) coordinate (O)
(-\c,0) coordinate (F_1)
(\c,0) coordinate (F_2)
(-\a,0) coordinate (A_1)
(\a,0) coordinate (A_2)
(0,-2) coordinate (B_1)
(0,2) coordinate (B_2)
(\xm,\ym) coordinate (M)
;
\draw (O)  let\p1=($(O)-(A_1)$) in circle ({veclen(\x1,\y1});
\draw (3,2)--(3,-2)--(-3,-2)--(-3,2)--(3,2) ;
\draw[fill=black] (A_1) circle(1pt);
\draw[fill=black] (A_2) circle(1pt);
\draw[fill=black] (B_1) circle(1pt);
\draw[fill=black] (B_2) circle(1pt);
\draw[fill=black] (O) circle(1pt);
%	\draw[dashed] (F_1)--(M)--(F_2)--(F_1);
\end{tikzpicture}}
\shortans[]{$21{,}5$}
\loigiai{\immini{Chọn hệ tọa độ $Oxy$ như hình vẽ.\\
Khi đó đường tròn có tâm $O$ và có độ dài đường kính bằng độ dài một cạnh nhỏ của hình chữ nhật là $ 16$ cm. Đường tròn này tiếp xúc với hypebol tại đỉnh của hypebol là $A_1(-8; 0)$ và $ A_2(8 ; 0)$.\\
Mặt khác ta thấy hypebol đi qua đỉnh $E(12;8)$ của hình chữ nhật.\\
Gọi dạng chính tắc của hypebol có phương trình $\break \dfrac{x^2}{a^2}-\dfrac{y^2}{b^2}=1$.\\
Vì hypebol có đỉnh $A_2(8;0)$ nên $a=8$.\\
Vì hypebol đi qua $E(12;8)$ nên $\dfrac{12^2}{a^2}-\dfrac{8^2}{b^2}=1$.
}{\begin{tikzpicture}[scale=0.8, font=\footnotesize, line join=round, line cap=round, >=stealth]
\def\a{2};
\def\b{1.788854382};
\pgfmathsetmacro\c{sqrt((\a)^2+(\b)^2)};
\draw[name path=h1,smooth, samples=200, domain=\a:3] plot (\x,{((\b)/(\a))*sqrt((\x)^2-(\a)^2)});
\draw[yscale=-1,name path=h2,smooth, samples=200, domain=\a:3] plot (\x,{((\b)/(\a))*sqrt((\x)^2-(\a)^2)});
\draw[xscale=-1,name path=h2,smooth, samples=200, domain=\a:3] plot (\x,{((\b)/(\a))*sqrt((\x)^2-(\a)^2)});
\draw[xscale=-1,yscale=-1,name path=h2,smooth, samples=200, domain=\a:3] plot (\x,{((\b)/(\a))*sqrt((\x)^2-(\a)^2)});
\def\xm{2};
\pgfmathsetmacro\ym{((\b)/(\a))*sqrt((\xm)^2-(\a)^2)};
\path
(0,0) coordinate (O)
(-\c,0) coordinate (F_1)
(\c,0) coordinate (F_2)
(-\a,0) coordinate (A_1)
(\a,0) coordinate (A_2)
(0,-2) coordinate (B_1)
(0,2) coordinate (B_2)
(\xm,\ym) coordinate (M)
(-\a,0) coordinate (A)
(\a,0) coordinate (B)
(0,-2) coordinate (C)
(0,2) coordinate (D)
(3,2) coordinate (E)
;
\draw (O)  let\p1=($(O)-(A_1)$) in circle ({veclen(\x1,\y1});
\draw (3,2)--(3,-2)--(-3,-2)--(-3,2)--(3,2) ;
\draw[fill=black] (A_1) circle(1pt);
\draw[fill=black] (A_2) circle(1pt);
\draw[fill=black] (B_1) circle(1pt);
\draw[fill=black] (B_2) circle(1pt);
\draw[fill=black] (O) circle(1pt);
%	\draw[dashed] (F_1)--(M)--(F_2)--(F_1);
\draw[->] (-4,0)--(4.5,0) node[below] { $x$};
\draw[->] (0,-2.5)--(0,3.5) node[left] { $y$};
\draw (0,0) node [below left] { $O$};
\foreach \x/\g in {A_1/-40,A_2/60,C/-50,E/0,D/50}
\fill[black] 	(\x) circle (1pt)($(\g:4mm)+(\x)$) node {$\x$};
\end{tikzpicture}}\noindent
Khi đó $\dfrac{12^2}{8^2}-\dfrac{8^2}{b^2}=1\Leftrightarrow \dfrac{5}{4}=\dfrac{8^2}{b^2}\Leftrightarrow b^2=\dfrac{256}{5}$.\\
Do đó $c=\sqrt{a^2+b^2}=\sqrt{64+\dfrac{256}{5}}=\dfrac{24\sqrt{5}}{5}\approx10{,}7$.
\\ Vậy tiêu cự của hypebol cần tìm là $2c\approx 21{,}5$.		}
\end{ex}

\begin{ex}%[0H9V4-2]
	Trong mặt phẳng tọa độ $Oxy$, cho điểm $A(-1;1)$, $B(0;-2)$, $C(0;2)$. Đường tròn ngoại tiếp tam giác $ ABC $ có phương trình là $ (x-a)^2+(y-b)^2=R^2 $. Tính $ S=a+b+R^2 $.
	\shortans{$6$}
	\loigiai{
		Giả sử tâm đường tròn là điểm $I(a;b)$.\\
		Ta có \allowdisplaybreaks
		\begin{eqnarray*}
			& &IA=IB=IC\Leftrightarrow IA^2=IB^2=IC^2\\
			&\Leftrightarrow&\heva{& (-1-a)^2+(1-b)^2=(0-a)^2+(-2-b)^2\\ & (0-a)^2+(-2-b)^2=(0-a)^2+(2-b)^2}\\
			&\Leftrightarrow& \heva{& a^2+b^2+2a-2b+2=a^2+b^2+4b+4\\ & a^2+b^2+4b+4=a^2+b^2-4b+4}\\
			&\Leftrightarrow& \heva{& 2a-2b=4b+2\\ & b=0}\Leftrightarrow\heva{& a=1\\ & b=0.}
		\end{eqnarray*}
		Đường tròn tâm $I(1;0)$, bán kính $R=IC=\sqrt{a^2+b^2-4b+4}=\sqrt{5}$.\\
		Suy ra phương trình đường tròn đi qua ba điểm $A$, $B$, $C$ có dạng $(x-1)^2+y^2=5$.\\
		Vậy $a=1$, $b=0$, $R^2=5$, $a+b+R^2=6$.
	}
\end{ex}

\Closesolutionfile{ans}

\TL
\begin{ex}
	Tìm số hạng chứa $x^4$ của khai triển nhị thức $\left(x^2-2y\right)^5$.
\end{ex}


\begin{ex}%[1 điểm]%[0H9V4-5]%[Dự án đề kiểm tra Toán 10 HKII NH23-24- Anh Hoàng Lê]%[THPT Chuyên Lê Hồng Phong - Tp HCM]
Trong mặt phẳng $Oxy$, cho đường tròn $(C)\colon x^2+y^2-2x+4y-4=0$ và đường thẳng $d\colon 2x+y+2 \sqrt{15}=0$. 
\begin{enumerate}
	\item Viết phương trình đường thẳng $\Delta$ đi qua tâm đường tròn $\left(C\right)$ và vuông góc với đường thẳng $d$.
	\item Có bao nhiêu điểm thuộc $d$ sao cho qua điểm đó có thể kẻ hai tiếp tuyến đến $(C)$ và góc giữa hai tiếp tuyến bằng $60^{\circ}$.
\end{enumerate}
\loigiai{
Gọi $M$ là điểm thỏa yêu cầu bài toán. Cho tiếp tuyến tại $M$ tiếp xúc với $(C)$ tại $A$ và $B$. Ta chia thành hai trường hợp.

\faStar~~\textbf{Trường hợp 1.} $\widehat{AMB}=60^{\circ}$.\\
Ta có $(C)\colon x^2+y^2-2x+4y-4=0$ hay $(x-1)^2+(y+2)^2=9$ suy ra $(C)$ có tâm là $I(1;-2)$ và bán kính $R=3$.\\
Do $M \in d$ nên ta gọi tọa độ của điểm $M$ là $\left(k;-2k-2 \sqrt{15}\right)$ suy ra
\[I M=|\overrightarrow{IM}|=\sqrt{\left(k-1\right)^2+\left(-2k-2 \sqrt{15}+2\right)^2}. \]
Mặt khác  $\widehat{IMA}=\widehat{IMB}=\dfrac{\widehat{AMB}}{2}=\dfrac{60^{\circ}}{2}=30^{\circ}$ và $I A \perp A M$ nên
\[I M=\dfrac{IA}{\sin \widehat{AMI}}=\dfrac{R}{\sin 30^{\circ}}= \dfrac{3}{\sin 30^{\circ}}=6.\] Suy ra
\[
\left(k-1\right)^2+\left(-2k-2 \sqrt{15}+2\right)^2=36 \,\, \text{hay} \,\,  5  k^2-2k\left(5-4 \sqrt{15}\right)+29-8 \sqrt{15}=0.
\]
Ta có $\Delta'=\left(5-4 \sqrt{15}\right)^2-5 \cdot \left(29-8 \sqrt{15}\right)=120>0$.\\
Vậy phương trình trên có hai nghiệm phân biệt hay có  $2$  điểm $M$ thỏa.

\faStar~~\textbf{Trường hợp 2.} $\widehat{AMB}=120^{\circ}$.\\
Tương tự như trên ta có $I M=\dfrac{3}{\sin 60^{\circ}}=2 \sqrt{3}$ nên ta có phương trình
\[
\left(k-1\right)^2+\left(-2k-2 \sqrt{15}+2\right)^2=12 \,\, \text{hay} \,\,  5k^2-2k\left(5-4 \sqrt{15}\right)+53-8 \sqrt{15}=0.
\]
Ta có $\Delta'=\left(5-4 \sqrt{15}\right)^2-5 \cdot \left(53-8 \sqrt{15}\right)=0$ hay phương trình trên có nghiệm kép.\\
Vậy có tất cả  $3$  điểm $M$ thỏa đề bài.
}
\end{ex}

\begin{ex}
	Một nhóm gồm $6$ học sinh nam và $6$ học sinh nữ. Xếp ngẫu nhiên $12$ học sinh này thành một vòng tròn. Tính xác suất để giữa mỗi hai học sinh nam hoặc có đúng $2$ học sinh nữ hoặc không có học sinh nữ nào.
\end{ex}
\Closesolutionfile{ansbook}
